\section{Expanded reasoning: mechanics, electronics and motion algorithms}
\label{sec:expand-hardware}
\subsection{Pulley travel, speed and torque: assumptions before arithmetic}
For belt pitch $p=2$ mm and $N_t=20$ teeth, travel per revolution is
$L_{\rm rev}=N_tp=40$ mm, and pitch radius
$r=L_{\rm rev}/(2\pi)=6.3662$ mm. This is the belt pitch radius, not
necessarily the measured outer pulley radius. At 4000 mm/min, motor speed
is $4000/40=100$ rev/min. The report rounded a cited 3600 mm/min handwriting
speed up to 4000, an increase of about 11.1\%; this is a design allowance,
not an experimentally recovered maximum gantry speed~\cite{finalreport2026}.

Torque is the moment of force; for tangential belt force $F$, $T=rF$.
The general vector definition is $\boldsymbol\tau=\mathbf r\times\mathbf F$
and magnitude $rF\sin\theta$~\cite{openstaxrotation}. Assuming steady
horizontal motion and a Coulomb-style equivalent friction $F_f=\mu mg$,
the report uses $\mu=0.004,m=0.950$ kg and $g=9.81$ m/s$^2$:
\[
F_f=0.037278\ {\rm N},\qquad
T_f=0.0063662(0.037278)=0.0002373\ {\rm N\,m}
       =0.02373\ {\rm N\,cm}.
\]
The conversion multiplies N\,m by 100 to obtain N\,cm. The small answer
follows from a small assumed friction coefficient and short radius. It does
not establish the actual rolling resistance of this assembly.

A fuller horizontal acceleration model is
\[
T_{\rm req}\simeq {r\over\eta}(ma+F_f+F_{\rm pen}+F_{\rm other})
              +J_{\rm rot}{a\over r},
\]
where $\eta$ is transmission efficiency, $F_{\rm pen}$ effective drag,
$F_{\rm other}$ other opposing forces and $J_{\rm rot}$ rotary inertia.
Terms must be referred to the same shaft. Constructed example: at
$a=0.4$ m/s$^2$, translational acceleration adds
$rma=0.002419$ N\,m $=0.2419$ N\,cm before other terms, over ten times
the report's steady-friction estimate. This illustrates an omitted load,
not a newly measured torque or a motor redesign.

The report interpolates curve points $(60,7)$ and $(120,21)$ in RPM and
N\,cm: $T(100)=7+(100-60)(21-7)/(120-60)=16.33$ N\,cm.
Linear interpolation assumes the correct curve and operating conditions.
The report literally leaves the curve citation as ``SOURCE''. The unusual
increasing pair cannot be independently authenticated here. Do not turn
16.33 into an established available torque or use its ratio to 0.02373 as
a validated safety factor. Pull-out torque at speed, holding torque at rest
and pull-in/start capability are different. Supply voltage, phase current,
driver mode and acceleration affect the comparison.

\subsection{Resolution, quantisation and why they are not accuracy}
For $N_s=200$ full steps/rev and microstep factor $m_s=16$,
$K=N_sm_s/L_{\rm rev}=80$ commanded steps/mm, spacing
$\Delta=1/K=0.0125$ mm. Full-step angle is $360/200=1.8$ degrees.
Nearest-grid rounding has per-axis error at most $\Delta/2=0.00625$ mm.
With equal axes, endpoint Euclidean rounding error is at most
$\sqrt2\Delta/2=0.00884$ mm. If two independently rounded endpoints define
a length, their relative position can differ by up to twice the endpoint
bound. Truncation instead of nearest rounding has a different bound.

These are ideal coordinate quantisation bounds. Microstep current ratios do
not ensure exactly equal loaded rotor angle increments. Belt compliance,
backlash, missed steps, pen flex and calibration error are separate. An
open-loop counter records requested position; it does not observe actual
position. Consequently the 0.5 mm physical length requirement requires
measurement even though the command lattice is much finer.

\subsection{A4988 current, reference voltage and every driver signal}
The manufacturer gives $I_{\rm trip,max}=V_{\rm REF}/(8R_S)$.
With confirmed $R_S=0.1\,\Omega$ and desired peak limit 1.68 A,
$V_{\rm REF}=8(0.1)(1.68)=1.344$ V~\cite{a4988datasheet}.
The factor 8 belongs to the driver's internal current-sense scale. It is
not the microstep factor, number of coils or supply voltage. For a different
carrier with $R_S=0.068\,\Omega$, the same limit would require 0.914 V.
Pololu explicitly warns that carrier sense resistance varies~\cite{pololu4988}.
Thus ``R10 means 0.1 ohm'' is a component-identification assumption requiring
the actual board marking, not a universal carrier value.

This equation sets an electrical current limit; it does not by itself prove
a thermally safe operating point. The final report's main text claims an
adjustment, whereas its connection table says actual $V_{\rm REF}$ remains
unmeasured. This edition retains that as unresolved provenance. The external
supply's displayed current is not coil current: chopping and energy storage
make input and winding currents different quantities. Full-step two-coil
current also differs from the peak microstep limit.

\begin{description}[style=nextline]
\item[STEP and DIR.] STEP's active edge advances the translator; DIR chooses
sequence direction. Logical inversion maps positive software travel to the
installed motor orientation. The datasheet lists minimum STEP high and low
times of 1 microsecond, and DIR setup/hold of 200 ns. These are electrical
minimums, not guarantees of Linux scheduling precision.
\item[MS1, MS2, MS3.] Select microstep mode; all high selects 1/16 in the
manufacturer's truth table. The source assumes this setting; it does not
read back physical straps.
\item[ENABLE, RESET, SLEEP.] Enable controls bridge output state, reset returns
the translator to a defined home state, and sleep reduces activity. Their
installed wiring is not fully documented in the final report. Do not infer
GPIO assignments merely from the driver's pin list.
\item[VDD and GND.] Logic supply and common signal reference, distinct from
motor supply. A shared reference is required for meaningful logic levels.
\item[VMOT / VBB and winding outputs.] Motor supply feeds the bridge;
1A/1B and 2A/2B connect to winding pairs. They carry switched winding
current, not GPIO STEP waveforms. Reversing a winding pair changes orientation.
\item[VREF and SENSE.] Reference and current-sense network establish the
regulated limit. They are analog/current-control nodes, not ordinary
logic commands. Carrier capacitors and charge-pump components have support
roles described in the datasheet; no undocumented installed values are assumed.
\end{description}

\subsection{BC337 low-side switch, base resistor and switch pull-down}
In a low-side NPN arrangement, motor positive goes to its supply, motor
negative to collector, emitter to common ground, and GPIO through a base
resistor. GPIO supplies base current, not motor current. The approximate
base relation is $I_B=(V_{\rm GPIO}-V_{BE})/R_B$~\cite{rohmnpn}.
Assuming 3.3 V, $V_{BE}=0.7$ V and the report's stated 4 mA limit,
$R_{B,\min}=(3.3-0.7)/0.004=650\,\Omega$.
The limit is a \emph{minimum resistance}; lowering it raises GPIO current.
At $R_B=1$ k$\Omega$, nominal $I_B=2.6$ mA and resistor dissipation is
$I_B^2R_B=6.76$ mW. Output-high droop and actual saturation $V_{BE}$
change that estimate.

The transistor's active-region DC gain is not a guaranteed saturation gain.
At a forced gain of ten, 2.6 mA base current corresponds to a 26 mA load
check. The report lists motor rated current 30 mA and stall current 230 mA;
these require forced gains 11.5 and 88.5 respectively at that nominal base
current. This arithmetic alone therefore does not establish saturation at
either load. The onsemi saturation-test conditions should be consulted
directly; automated retrieval was blocked~\cite{bc337data}. No circuit
change or pressure fix is performed here.

An inductive motor stores $E=\tfrac12LI^2$. Interrupting current produces
$v=L\,di/dt$, explaining why turn-off behaviour differs from a resistor.
The report records no fitted flyback diode; this manual does not invent one.
The physical transient and transistor dissipation cannot be quantified without
motor inductance, current and switching measurements.

For an active-high switch with 10 k$\Omega$ pull-down, an open input is
held near ground; closure connects it to the high rail. At 3.3 V, closed
pull-down current is 0.33 mA and resistor power 1.089 mW.
For input leakage $I_{\rm leak}$, open-node error is approximately
$I_{\rm leak}R$; cable/input capacitance gives time constant $RC$.
A weaker pull-down saves current but increases susceptibility to leakage
and coupled noise. Ten kilohms is a common engineering scale, not a value
proved optimal by these equations. A base-emitter pull-down has a different
purpose from a switch-input pull-down; the report does not fully establish
the placement of its stated base pull-down.

\subsection{All documented controller pins and supply connections}
Header number means physical connector position; line offset means a GPIO
line within the requested chip. They are not interchangeable~\cite{kernelgpio}.
The local source uses \texttt{/dev/gpiochip0}. The following mapping agrees
with report Table 17 and inspected source, not a new electrical pin test.
The manufacturer expansion reference was inaccessible to the automated
reader~\cite{odroidgpio}.

\begin{center}\small
\begin{tabular}{rrll}
\toprule Header & Line & Signal & Role \\\midrule
7&64& DIR1 & X direction output\\
11&68& STEP1 & X step output\\
12&81& DIR2 & Y direction output\\
13&69& STEP2 & Y step output\\
15&72& X minimum & Active-high switch input\\
16&65& X maximum & Active-high switch input\\
18&66& Y minimum & Active-high switch input\\
22&67& Y maximum & Active-high switch input\\
33&70& Pen switch & Raised-state feedback input\\
35&71& Pen motor & Transistor-base control output\\
\bottomrule
\end{tabular}
\end{center}
The report separately lists USB camera data/power, network communication,
ODROID DC input, external 24 V to driver motor supply, a dedicated 3.3 V
rail for logic/pen supply, and common grounds. It does not record exact
physical power/ground pins, complete return routing, actual adapter output
or total 3.3 V rail load. USB and Ethernet are interfaces, not spare GPIO
numbers. The quoted 4 mA GPIO capability is conditional because the report
leaves its specific source unresolved; it must not be inferred from the
stepper driver or BC337 datasheet.

\subsection{Pen state machine, dwell, overrun and debounce}
The pen motor turns one way and the cam alternates states. The conceptual
state machine is: check interlocks; read current state; if already at target,
do nothing; otherwise energise; poll until target or timeout; optionally
continue the pen-up seating overrun; de-energise; report success/failure.
The inspected \texttt{SetPen} compares feedback with the requested target,
not always ``up'' as one report pseudocode line suggests. A switch reporting
not-up does not measure nib contact force. The mechanism can satisfy a
binary switch while pen pressure remains unsuitable.

The 40 ms up-overrun, 2 ms polling sleep, 3 s timeout and 350 ms G-code
settling dwell serve different roles. The report attributes dwell to visual
tuning; it supplies no formal optimum for each. At an illustrative unloaded
150 RPM, a quarter-turn takes 0.1 s, but gearing load, cam contact and
switch position make this unsuitable as a measured timing prediction.
The same 350 ms dwell need not represent the same physical settling process
for raising and lowering.

Endstop confirmation requires three high reads, two gaps of 6 ms, hence
12 ms from first to last sample. Under an unrealistic independent spurious
high probability $q$, three samples would give $q^3$. Real contact bounce
is time-correlated, so this is not a reliability estimate. Debounce rejects
short chatter at a cost of latency; it cannot ensure switches are ``never''
triggered by noise. The pen feedback routine is distinct from the debounced
endstop routine. At nominal 7.8125 mm/s, 12 ms corresponds to 0.09375 mm
of travel \emph{if motion continued}; the inspected polling can pause pulses
during confirmation, so that multiplication is not a measured stopping distance.

\subsection{Calibration and G-code motion mapping}
For each axis, move to a switch, reverse, count $N$ pulses to the other
switch, use measured travel $L$ to obtain $K=N/L$ steps/mm, establish which
end is minimum/maximum, then retract 5 mm. With $L=200.6$ mm, the two 5 mm
end clearances give 190.6 mm usable span; the recorded Y travel 261.3 gives
251.3 mm. The 5 mm is a clearance choice, not an inferred switch precision.
Homing at 1.6 ms per nominal step gives
$1/(0.0016\cdot80)=7.8125$ mm/s before overhead.
The 200000-step bound is a failure guard, not measured travel.

Calibration estimates scale from an external length and reference contacts.
It does not calibrate squareness, backlash or flex. In $K=N/L$, errors in
$L$ propagate approximately as $\delta K/K\simeq-\delta L/L$ for fixed $N$.
Calibrating against the same length later used for acceptance can conceal
systematic error, so provenance of independent checks matters.

G21 selects millimetres; G90 absolute positions; G0/G1 specify linear targets
with feed handling; F is mm/min, so 900 becomes 15 mm/s and 2400 becomes
40 mm/s. G4 P is a dwell with this project's P convention in seconds.
M3/M5 are repurposed here to lower/raise the pen; in other controllers they
often refer to spindle control. M2 terminates. Modal settings persist until
changed. Parsing a command and updating a commanded counter does not verify
that the carriage reached its target.

\subsection{Bresenham movement and how to read its staircase}
For integer pulse displacements $(d_x,d_y)$, set
$a_x=|d_x|,a_y=|d_y|$ and direction signs. If $a_x\ge a_y$, initialise
$e=\lfloor a_x/2\rfloor$. At each X pulse subtract $a_y$; if $e<0$,
also pulse Y and add $a_x$. Swap axes for a steep line. This distributes
minor pulses along major pulses with integer arithmetic~\cite{bresenham1965}.

For $(5,2)$, errors after each complete iteration are $0,3,1,4,2$ and
the endpoint sequence is
\[
(1,0),\ (2,1),\ (3,1),\ (4,2),\ (5,2).
\]
The ideal Y coordinates are $0.4,0.8,1.2,1.6,2.0$; discrepancies are
$-0.4,0.2,-0.2,0.4,0$ steps. Signs reflect axes; the absolute deviations
are at most half a minor step in this example. The plot's staircase shows
the command lattice, not measured carriage vibration. An equal-axis move
pulses both axes per iteration. A zero-length move needs no pulses.

Geometry and timing are distinct. The inspected routine pulses major and
minor axes sequentially, sleeping for each high pulse. Approximately
$a_x+a_y$ pulse sleeps are therefore incurred. For nominal path length
$L$ and feed $v$, a naive $N_{\rm major}$-based delay can make diagonal
motion slower unless this sequential cost is accounted for. GPIO calls,
debounce and OS scheduling add overhead. Bresenham alone supplies neither
acceleration control nor hard-real-time timing.

\subsection{Bounds, joins and nominal timing budget}
For polyline vertices $p_i=(x_i,y_i)$, bounds are the coordinate minima and
maxima. Straight segments lie within their endpoint box, so all vertices
inside a rectangular workspace imply the interpolated straight segments do
too. This does not cover pen radius, calibration uncertainty, cam geometry
or curved interpolation extrema. The empty-trajectory zero box is a
convention, not a drawn dot at the origin.

Pen-down length is $L_d=\sum\|p_{i+1}-p_i\|$ within strokes; pen-up travel
is the distance between successive stroke endpoints. At feeds $F_d,F_u$
in mm/min, a nominal time is
$t_0=60L_d/F_d+60L_u/F_u+N_{\rm toggle}\tau$ plus any explicitly counted
dwells. For $L_d=300$ mm, $L_u=120$ mm, feeds 900/2400 and twenty 0.35 s
dwells, $t_0=20+3+7=30$ s. This is constructed ideal accounting, not an
ODROID prediction. Avoid double-counting a dwell already included in a
toggle term. Acceleration, sequential pulses, software overhead and
generation time explain why measured full-request timing differs.

The 0.6 mm endpoint-join threshold avoids a lift/drop pair by drawing a
short connector. It trades fewer toggles for a possible extra mark; at
$H=4$ mm it is 0.15 x-heights. It is not a Douglas--Peucker error bound.
Ring closure, decimation and stroke-order optimisation likewise change
different aspects of a trajectory. Their combined output must be considered
when interpreting physical lengths or geometric fidelity; none is tuned or
tested by this document-only expansion.
